\section{Technical Evaluation}

A document is only worth returning to if it still points at the code it describes, and able to keep up with the latest codebase faithfully~\cite{fluri2007code, wen2019large}. We tested that property directly on five open-source projects we rewound the repository to a commit in its past, built the document there, and then applied every later commit in the order its own developers wrote them, checking after each one whether every section still reached the code it was about. The correct answer for each commit comes from the commit itself and not from \sys{}. We recover where each definition ended up by matching names across the two versions of the files a commit touched, which is a signal \sys{} never uses, so agreement between the two is evidence rather than restatement. Five further projects were used to build and debug the harness and are excluded from the figures below, and the system was frozen before this run. Appendix~\ref{app:attachment} gives the protocol and the per-project results.

Across 615 commits and 26{,}250 attached regions of code, no section ever pointed at code that was no longer there. All 125 deletions were released. Of the 3.3 million times a section pointed into a file that a commit did not touch, none moved. Of the 74 definitions that moved to another file, one lost its section. Between a third and three quarters of commits were handled by mechanical rules with no model involved, so the model is consulted only to place code that has just been written, which is what keeps the cost bounded.

A whole-repository restructure defeats this. When one project renamed its top-level directory, 440 files and 6{,}556 definitions moved without a line of their contents changing, and 6{,}275 of 7{,}209 sections lost their code in that single commit. The step that re-reads changed files scales badly with how much moves at once and reported almost none of the change to the rest of the system. Ordinary development stays well inside what works, including the refactoring present in the five projects above. We report this as an open defect rather than a limitation of the approach, because the mechanism that should have absorbed the rename is deterministic and needs no model. Appendix~\ref{app:rename} records the measurements and a reproduction.

Staying attached should also make the document more useful to an agent, and that is measurable. We kept one copy of the document maintained across 200 commits and froze a second copy at the starting commit, then asked an agent to find the code behind 25 later change requests that neither copy had seen. The agent given the maintained copy identified the right files more often, recovering 40\% of them against 33\% with the frozen copy and 29\% with no document at all, and it used a quarter fewer tokens doing so. A stale reference sends an agent to files that no longer exist and it pays for the detour. On a project whose code barely moves the three conditions were indistinguishable, so maintenance earns its keep where code moves, which is where documents rot. Appendix~\ref{app:drift} reports the task construction and the per-condition figures.
